\documentclass{article}
\usepackage{graphicx} % Required for inserting images
\usepackage{fontspec}
\usepackage{polyglossia}
\setdefaultlanguage{greek}
\setotherlanguage{english}
\setmainfont{GFS Artemisia}
\title{Autonomous exploration and mapping}
\author{Πλιάκου Ελένη, 1092810 
    \\ Παπαχατζάκη Ανδριάνα, 1100996 }
\date{Σεπτέμβριος 2026}

\begin{document}

\maketitle

\section{Εισαγωγή}
Το πρότζεκτ αφορά στη αυτόνομη χαρτοχράφηση και εξερεύνηση ενός αρχικά άγνωστου χώρου, με δεδομένες διαστάσεις, συμπεριλαμβανομένης της αποφυγής εμποδίων, με ενα crazyflie 2.1. Προκειμένου να υλοποιηθούν αυτά τα ζητούμενα, χρησιμοποιήθηκαν οι ακόλουθοι αισθητήρες:
\begin{itemize}
    \item Lighthouse, που με τους φωτοαισθητήρες, λαμβάνει την υπερυθρη ακτινοβολία απο τους σταθμούς βάσης, ωστε να εντοπίσει την ακριβή του θέση στον χώρο
    \item Multirange, που με τους πέντε αισθητήρες του εντοπίζει την απόσταση από τα εμπόδια μπροστα, πίσω , δεξιά, αριστερά και πάνω%στις τεσσερις κατευθύνσεις
    \item Flow deck, που με έναν αισθητήρα υπολογίζει την απόσταση από το έδαφος, και έναν άλλο μπορεί να εκτιμήσει την οριζόντια κίνηση του 
\end{itemize}
Επομένως το προτζεκτ χωρίζεται σε δύο μέρη, το πρώτο που είναι η υλοποίηση του με lighthouse και multirange, και το δεύτερο που είναι η υλοποίηση με flowdeck και multirange.
\section{Υλοποίηση}
Προκειμένου να προχωρήσουμε στην υλοποίηση πρέπει να καταγράψουμε όλες τις διεργασίες, αυτές είναι:
\begin{itemize}
    \item Αυτόνομη απογείωση
    \item Εξερεύνηση του άγνωστου περιβάλλοντος
    \item Εκτίμηση της κίνησης και της θέσης, και η αναγνώριση εμποδίων 
    \item Δημιουργία ενός συνεχώς ανανεούμενου χάρτη που αναπαριστά το περιβάλλον, και τον διαχωρισμό των περιοχών που έχουν επισκεφθεί
    \item αναγνώριση των περιοχών σε ανεξερεύνητες περιοχές
    \item Υπολογισμός τροχιών μέσω διαθέσιμων μονοπατιών
    \item Συνεχης εξερεύνηση μέχρι την πλήρη κάλυψη του
    \item Επιστροφή στο σημείο εκκίνησης και προσγείωση
\end{itemize}

\end{document}

